\documentclass[11pt,a4paper]{article}

% ============ PAQUETES ============
\usepackage[utf8]{inputenc}
\usepackage[T1]{fontenc}
\usepackage[spanish]{babel}
\usepackage{geometry}
\geometry{margin=2.3cm}
\usepackage{graphicx}
\usepackage{tikz}
\usetikzlibrary{shapes,arrows.meta,positioning,fit,backgrounds,calc}
\usepackage{booktabs}
\usepackage{longtable}
\usepackage{array}
\usepackage{xcolor}
\usepackage{enumitem}
\usepackage{hyperref}
\hypersetup{colorlinks=true, linkcolor=blue, urlcolor=blue, citecolor=blue}
\usepackage{fancyhdr}
\usepackage{titlesec}
\usepackage{amssymb}

\titleformat{\section}{\normalfont\Large\bfseries}{P\thesection}{0.6em}{}
\renewcommand{\thesection}{\arabic{section}}

\pagestyle{fancy}
\fancyhf{}
\fancyhead[L]{Prueba Práctica -- Unidad IV -- ISR-401}
\fancyhead[R]{\thepage}
\renewcommand{\headrulewidth}{0.4pt}

\newcolumntype{L}[1]{>{\raggedright\arraybackslash}p{#1}}
\newcolumntype{C}[1]{>{\centering\arraybackslash}p{#1}}

% ============ DATOS DEL ESTUDIANTE ============
\newcommand{\nombreAlumno}{Danela Arteaga Alava}
\newcommand{\cedulaAlumno}{0941375016}
\newcommand{\curso}{4to ``Software A''}
\newcommand{\fechaEntrega}{11/08/2026}
\newcommand{\urlRepo}{https://github.com/darteagaa-boop/prueba-practica-isr401-SEMANA15}
% \nombreAlumno, \curso, \fechaEntrega y \urlRepo deben ajustarse a los datos reales
% antes de la entrega final.

\begin{document}

% ==================================================================
% CARÁTULA
% ==================================================================
\begin{titlepage}
\centering
{\large UNIVERSIDAD TÉCNICA ESTATAL DE QUEVEDO\par}
{\large Facultad de Ciencias de la Ingeniería\par}
{\large Carrera de Ingeniería de Software\par}
\vspace{1.2cm}
{\Huge\bfseries Prueba Práctica -- Unidad IV\par}
\vspace{0.4cm}
{\Large Validación, Gestión, Herramientas y Estándares de Requisitos\par}
\vspace{0.3cm}
{\large Asignatura: Ingeniería de Requisitos (ISR-401)\par}
{\large Docente: Ing. Gleiston Guerrero, Mg.\par}
\vspace{0.7cm}

\begin{tabular}{ll}
\textbf{Estudiante:} & \nombreAlumno \\[4pt]
\textbf{Cédula:} & \cedulaAlumno \\[4pt]
\textbf{Curso:} & \curso \\[4pt]
\textbf{Fecha:} & \fechaEntrega \\[4pt]
\textbf{Caso:} & Sistema de Gestión de Pedidos \\
\end{tabular}

\vspace{0.7cm}
\textbf{\large URL del repositorio (línea única):}\\[6pt]
{\large\url{\urlRepo}}

\vspace{0.6cm}
\textbf{\normalsize Captura del resumen del cuestionario (SGA)}\\[4pt]
\fbox{\includegraphics[width=0.6\linewidth]{figuras/resumen_sga.png}}

\vspace{0.5cm}
\textbf{\normalsize Captura de la evaluación / intento (SGA)}\\[4pt]
\fbox{\includegraphics[width=0.6\linewidth]{figuras/intento_sga.png}}

\vfill
\end{titlepage}

\tableofcontents
\newpage

% ==================================================================
\section{Modelo de datos -- Diagrama de clases UML}
\label{sec:p1}

Diagrama de clases UML del dominio del \emph{Sistema de Gestión de Pedidos}, con las
clases \textbf{Cliente}, \textbf{Pedido}, \textbf{LíneaPedido} y \textbf{Producto}, sus
atributos tipados, una operación relevante por clase y las multiplicidades en cada
asociación. Los identificadores de cada clase se muestran subrayados.

\begin{figure}[h!]
\centering
\resizebox{\textwidth}{!}{%
\begin{tikzpicture}[
  class/.style={rectangle, draw=black, thick, inner sep=0pt},
  >=Stealth
]

% ---- Cliente ----
\node[class] (cliente) at (0,0) {%
  \begin{tabular}{@{\ }l@{\ }}
  \textbf{Cliente} \\
  \hline
  \underline{cedulaRUC : string} \\
  nombre : string \\
  correo : string \\
  \hline
  + realizarPedido() \\
  \end{tabular}%
};

% ---- Pedido ----
\node[class, right=3.6cm of cliente] (pedido) {%
  \begin{tabular}{@{\ }l@{\ }}
  \textbf{Pedido} \\
  \hline
  \underline{idPedido : string} \\
  fecha : date \\
  total : decimal \\
  estado : string \\
  \hline
  + confirmarPedido() \\
  \end{tabular}%
};

% ---- LineaPedido ----
\node[class, right=3.6cm of pedido] (linea) {%
  \begin{tabular}{@{\ }l@{\ }}
  \textbf{LíneaPedido} \\
  \hline
  cantidad : integer \\
  subtotal : decimal \\
  \hline
  + calcularSubtotal() \\
  \end{tabular}%
};

% ---- Producto ----
\node[class, below=2.4cm of linea] (producto) {%
  \begin{tabular}{@{\ }l@{\ }}
  \textbf{Producto} \\
  \hline
  \underline{código : string} \\
  descripcion : string \\
  precio : decimal \\
  stock : integer \\
  \hline
  + verificarStockProducto() \\
  \end{tabular}%
};

\draw[-] (cliente) -- (pedido);
\node[font=\scriptsize] at ($(cliente.east)+(0.5,0.35)$) {1};
\node[font=\scriptsize] at ($(pedido.west)+(-0.5,0.35)$) {0..n};

\draw[-] (pedido) -- (linea);
\node[font=\scriptsize] at ($(pedido.east)+(0.5,0.35)$) {1};
\node[font=\scriptsize] at ($(linea.west)+(-0.5,0.35)$) {1..n};

\draw[-] (linea) -- (producto);
\node[font=\scriptsize] at ($(linea.south)+(-0.4,-0.25)$) {1..n};
\node[font=\scriptsize] at ($(producto.north)+(-0.4,0.25)$) {1};

\end{tikzpicture}%
}
\caption{Diagrama de clases UML del dominio (P1).}
\label{fig:p1}
\end{figure}

\textbf{Multiplicidades:} un \emph{Cliente} realiza 0..n \emph{Pedidos}; un \emph{Pedido}
contiene 1..n \emph{LíneaPedido}; cada \emph{LíneaPedido} referencia exactamente 1
\emph{Producto} y un \emph{Producto} puede aparecer en 1..n \emph{LíneaPedido}.

\newpage
% ==================================================================
\section{Modelo funcional -- Diagrama de actividades UML}
\label{sec:p2}

Proceso \textbf{``Registrar pedido''}, modelado con carriles de responsabilidad
(\emph{Cliente} / \emph{Sistema}), nodo inicial, una decisión (\textit{¿Hay stock?}),
los dos caminos alternativos (confirmar pedido / notificar faltante), su convergencia
y el nodo final.

\begin{figure}[h!]
\centering
\begin{tikzpicture}[
  act/.style={rectangle, rounded corners, draw=black, thick, align=center, minimum width=3.1cm, minimum height=0.9cm, font=\small},
  dec/.style={diamond, draw=black, thick, align=center, font=\small, inner sep=1pt},
  >=Stealth
]

% Carriles
\draw[thick] (-2.2,0.6) rectangle (7.6,-9.4);
\draw[thick] (2.7,0.6) -- (2.7,-9.4);
\node[font=\bfseries] at (0.25,0.2) {Cliente};
\node[font=\bfseries] at (5.15,0.2) {Sistema};

% Nodo inicial
\node[circle, draw, fill=black, minimum size=0.35cm] (ini) at (0.25,-0.6) {};

% Registrar datos
\node[act] (reg) at (0.25,-1.9) {Registrar datos\\ del pedido};

% Verificar stock
\node[act] (ver) at (5.15,-1.9) {Verificar stock};

% Decisión
\node[dec] (dec) at (5.15,-3.7) {¿Hay\\ stock?};

% Confirmar pedido
\node[act] (conf) at (0.25,-5.6) {Confirmar pedido};

% Notificar faltante
\node[act] (not) at (5.15,-5.6) {Notificar producto\\ faltante};

% Nodo final
\node[circle, draw, fill=black, minimum size=0.35cm] (fin) at (0.25,-8.6) {};
\node[circle, draw, minimum size=0.55cm] at (0.25,-8.6) {};

\draw[->] (ini) -- (reg);
\draw[->] (reg) -- (ver);
\draw[->] (ver) -- (dec);
\draw[->] (dec) -- node[above,font=\scriptsize]{Sí} (conf);
\draw[->] (dec) -- node[right,font=\scriptsize]{No} (not);
\draw[->] (conf) -- (fin);
\draw[->] (not.south) -- ++(0,-1.5) -| (fin);

\end{tikzpicture}
\caption{Diagrama de actividades UML del proceso ``Registrar pedido'' (P2).}
\label{fig:p2}
\end{figure}

\newpage
% ==================================================================
\section{Modelo de comportamiento -- Máquina de estados UML}
\label{sec:p3}

Máquina de estados del ciclo de vida de un \emph{Pedido}. Se factoriza mediante un
superestado (\textit{Pedido Activo}) que agrupa los estados \textit{Registrado},
\textit{Confirmado}, \textit{Enviado} y \textit{Entregado}; la transición
\texttt{anularPedido} sale del superestado hacia \textit{Cancelado} (con guarda
\textit{[no entregado]}), evitando repetirla en cada subestado. \textit{Entregado} y
\textit{Cancelado} son estados finales de aceptación del ciclo de vida.

\begin{figure}[h!]
\centering
\begin{tikzpicture}[
  st/.style={rectangle, rounded corners, draw=black, thick, minimum width=2.6cm, minimum height=0.8cm, font=\small},
  >=Stealth
]

\node[circle, draw, fill=black, minimum size=0.3cm] (ini) at (0,0) {};

\begin{scope}[on background layer]
\node[draw, thick, rounded corners, fit={(-1.6,-0.9) (1.6,-6.6)}, inner sep=8pt, label={[font=\small\bfseries]above:Pedido Activo}] (super) {};
\end{scope}

\node[st] (reg) at (0,-1.3) {Registrado};
\node[st] (conf) at (0,-2.9) {Confirmado};
\node[st] (env) at (0,-4.5) {Enviado};
\node[st] (ent) at (0,-6.1) {Entregado};

\node[st] (can) at (0,-8.4) {Cancelado};
\node[circle, draw, minimum size=0.5cm] (finC) at (0,-9.8) {};
\node[circle, draw, fill=black, minimum size=0.3cm] at (0,-9.8) {};

\draw[->] (ini) -- (reg);
\draw[->] (reg) -- node[right,font=\scriptsize,align=left]{confirmarPedido\\ (stock mayor a 0)} (conf);
\draw[->] (conf) -- node[right,font=\scriptsize]{despacharPedido} (env);
\draw[->] (env) -- node[right,font=\scriptsize]{entregarPedido} (ent);

\draw[->] (super.east) -- ++(1.6,0) |- node[pos=0.25, right, font=\scriptsize, align=left]{anularPedido\\ (no entregado)} (can.east);

\draw[->] (can) -- (finC);

\end{tikzpicture}
\caption{Máquina de estados UML del ciclo de vida de un Pedido (P3).}
\label{fig:p3}
\end{figure}

\newpage
% ==================================================================
\section{Consistencia entre las tres perspectivas}
\label{sec:p4}

Tabla de verificación cruzada entre el evento (P3), la acción de actividad (P2) y el
elemento de datos que modifica (P1).

\begin{table}[h!]
\centering
\begin{tabular}{lll}
\toprule
\textbf{Evento (P3)} & \textbf{Acción (P2)} & \textbf{Elemento (P1)} \\
\midrule
confirmarPedido & Confirmar Pedido & Pedido.estado \\
despacharPedido & Despachar Pedido & Pedido.estado \\
entregarPedido  & Entregar Pedido  & Pedido.estado \\
anularPedido    & Anular Pedido    & Pedido.estado \\
verificarStock  & Verificar Stock  & Producto.stock \\
\bottomrule
\end{tabular}
\caption{Tabla de verificación cruzada P1--P2--P3.}
\label{tab:p4}
\end{table}

\subsection*{Inconsistencia detectada y corrección}

Se detecta que el evento \texttt{despacharPedido} existe en la máquina de estados
(P3), pero \textbf{no} tiene una acción equivalente en el diagrama de actividades
(P2), el cual solo modela ``Registrar pedido'' hasta su confirmación o el aviso de
faltante.

\textbf{Corrección aplicada:} se extiende conceptualmente el flujo del proceso de
negocio, añadiendo -- a continuación de ``Confirmar pedido'' -- las actividades
\textit{Despachar pedido} y \textit{Entregar pedido} en el carril \textit{Sistema},
de modo que el diagrama de actividades completo cubra todo el ciclo de vida
representado en la máquina de estados y quede así trazado en la Tabla~\ref{tab:p4}.

\newpage
% ==================================================================
\section{Especificación de requisitos con esquema de atributos}
\label{sec:p5}

Se especifican 2 requisitos funcionales y 2 no funcionales. Cada requisito no
funcional referencia una característica de \textbf{ISO/IEC 25010:2023} con métrica y
umbral.

\begin{table}[h!]
\centering\small
\begin{tabular}{ll}
\toprule
\multicolumn{2}{l}{\textbf{RF-01 -- El sistema debe registrar pedidos de clientes}} \\
\midrule
Prioridad & Alta \\
Estado & Aprobado \\
Fuente & Cliente \\
Estabilidad & Alta \\
Versión & 1.0 \\
Verificación & Prueba funcional \\
\bottomrule
\end{tabular}
\hfill
\begin{tabular}{ll}
\toprule
\multicolumn{2}{l}{\textbf{RF-02 -- El sistema debe verificar stock antes}} \\
\multicolumn{2}{l}{\textbf{de confirmar un pedido}} \\
\midrule
Prioridad & Alta \\
Estado & Aprobado \\
Fuente & Empresa \\
Estabilidad & Alta \\
Versión & 1.0 \\
Verificación & Prueba funcional \\
\bottomrule
\end{tabular}

\vspace{0.6cm}

\begin{tabular}{ll}
\toprule
\multicolumn{2}{l}{\textbf{RNF-01 -- Eficiencia de desempeño (ISO/IEC 25010:2023)}} \\
\multicolumn{2}{l}{El sistema deberá registrar pedidos en menos de 3 segundos} \\
\midrule
Métrica & Tiempo de respuesta \\
Umbral & $<$ 3 segundos \\
Prioridad & Alta \\
Estado & Aprobado \\
Fuente & Empresa \\
Estabilidad & Media \\
Versión & 1.0 \\
Verificación & Prueba de rendimiento \\
\bottomrule
\end{tabular}
\hfill
\begin{tabular}{ll}
\toprule
\multicolumn{2}{l}{\textbf{RNF-02 -- Seguridad (ISO/IEC 25010:2023)}} \\
\multicolumn{2}{l}{Los datos personales deberán almacenarse cifrados} \\
\midrule
Métrica & Porcentaje de datos cifrados \\
Umbral & 100\% \\
Prioridad & Alta \\
Estado & Aprobado \\
Fuente & Legislación \\
Estabilidad & Alta \\
Versión & 1.0 \\
Verificación & Auditoría de seguridad \\
\bottomrule
\end{tabular}
\caption{Esquema de atributos de los requisitos RF-01, RF-02, RNF-01 y RNF-02.}
\end{table}

\newpage
% ==================================================================
\section{Priorización MoSCoW}
\label{sec:p6}

\begin{table}[h!]
\centering
\begin{tabular}{lll}
\toprule
\textbf{Requisito} & \textbf{Categoría} & \textbf{Justificación} \\
\midrule
RF-01  & Must   & Sin registro no existe pedido. \\
RF-02  & Must   & Evita ventas sin disponibilidad de stock. \\
RNF-01 & Should & Mejora la experiencia del usuario. \\
RNF-02 & Could  & Se implementa gradualmente, aunque es importante para la legislación. \\
\bottomrule
\end{tabular}
\caption{Priorización MoSCoW de los requisitos (P6).}
\end{table}

\newpage
% ==================================================================
\section{Validación por inspección (ISO/IEC/IEEE 29148)}
\label{sec:p7}

\subsection*{Defectos detectados}

\begin{table}[h!]
\centering
\begin{tabular}{ll}
\toprule
\textbf{Requisito} & \textbf{Defecto} \\
\midrule
RF-01  & No especifica qué datos registra. \\
RF-02  & No define qué ocurre cuando no hay stock. \\
RNF-01 & No indica la cantidad de usuarios concurrentes. \\
RNF-02 & No especifica el algoritmo de cifrado. \\
\bottomrule
\end{tabular}
\caption{Defectos detectados por inspección (identificación, sin corregir).}
\end{table}

\subsection*{Requisitos corregidos (retrabajo)}

\begin{itemize}[leftmargin=1.4cm]
\item \textbf{RF-01:} El sistema registrará pedidos con ID, fecha, cliente y total.
\item \textbf{RF-02:} El sistema verificará stock y notificará los faltantes antes
      de confirmar el pedido.
\item \textbf{RNF-01:} El sistema registrará pedidos en menos de 3 segundos para
      hasta 100 usuarios concurrentes.
\item \textbf{RNF-02:} El sistema almacenará los datos personales utilizando
      cifrado AES-256.
\end{itemize}

\newpage
% ==================================================================
\section{Pruebas de aceptación trazadas}
\label{sec:p8}

\begin{table}[h!]
\centering\small
\begin{tabular}{lllll}
\toprule
\textbf{Prueba} & \textbf{Requisito} & \textbf{Tipo} & \textbf{Entrada} & \textbf{Criterio de éxito} \\
\midrule
PA-01 & RF-01  & Funcional     & Pedido válido        & Pedido registrado \\
PA-02 & RF-02  & Funcional     & Producto sin stock    & Notificación mostrada \\
PA-03 & RNF-01 & No funcional  & 100 usuarios concurrentes & Tiempo de respuesta $<$ 3 s \\
PA-04 & RNF-02 & No funcional  & Auditoría de datos    & Datos almacenados cifrados \\
\bottomrule
\end{tabular}
\caption{Pruebas de aceptación trazadas a cada requisito (P8).}
\end{table}

\newpage
% ==================================================================
\section{Matriz de trazabilidad}
\label{sec:p9}

\begin{table}[h!]
\centering\small
\begin{tabular}{lllll l}
\toprule
\textbf{Req.} & \textbf{Fuente} & \textbf{P1 (Clase)} & \textbf{P2 (Actividad)} & \textbf{P3 (Estado)} & \textbf{Prueba} \\
\midrule
RF-01  & Cliente     & Pedido   & Registrar pedido      & Registrado & PA-01 \\
RF-02  & Empresa     & Producto & Verificar stock        & Confirmado & PA-02 \\
RNF-01 & Empresa     & Pedido   & Registrar / Confirmar pedido & Confirmado & PA-03 \\
RNF-02 & Legislación & Cliente  & Registrar datos del pedido & Confirmado & PA-04 \\
\bottomrule
\end{tabular}
\caption{Matriz de trazabilidad: fuente (traza hacia atrás) y modelos/pruebas (traza hacia adelante).}
\end{table}

\subsection*{Dependencia horizontal}
\textbf{RF-02} depende de \textbf{RF-01}: no es posible verificar el stock de un
pedido (RF-02) si el pedido aún no ha sido registrado (RF-01).

\subsection*{Requisitos huérfanos}
No existen requisitos huérfanos: los cuatro requisitos (RF-01, RF-02, RNF-01,
RNF-02) cuentan con al menos un modelo asociado (P1--P3) y con al menos una prueba
de aceptación asociada (P8).

\newpage
% ==================================================================
\section{Gestión del cambio y línea base}
\label{sec:p10}

\subsection*{Solicitud de cambio}
\textbf{SC-01:} Se solicita que el tiempo de registro de un pedido pase de 3
segundos a 2 segundos.

\subsection*{Análisis de impacto (según matriz de P9)}
La solicitud SC-01 afecta a:
\begin{itemize}[leftmargin=1.4cm]
\item RNF-01 (Eficiencia de desempeño).
\item PA-03 (prueba de aceptación de RNF-01).
\item El diagrama de actividades (P2), en el tiempo estimado de la actividad
      ``Registrar datos del pedido''.
\item La documentación de requisitos (ERS).
\end{itemize}

\subsection*{Decisión del CCB (Change Control Board)}
\textbf{Aprobada.} \checkmark

\subsection*{Actualización de versión}
\begin{center}
\begin{tabular}{ll}
\toprule
Requisito & Cambio de versión \\
\midrule
RNF-01 & v1.0 $\rightarrow$ v1.1 \\
\bottomrule
\end{tabular}
\end{center}

\textbf{Nuevo texto del requisito RNF-01 (v1.1):} ``El sistema registrará pedidos en
menos de 2 segundos para hasta 100 usuarios concurrentes.''

\subsection*{Línea base 1.1}
Se declara la siguiente línea base, con la configuración congelada y coherente en
versión 1.1:

\begin{itemize}[leftmargin=1.4cm]
\item ERS v1.1
\item Diagrama de clases v1.1
\item Diagrama de actividades v1.1
\item Máquina de estados v1.1
\item Matriz de trazabilidad v1.1
\item Pruebas de aceptación v1.1
\end{itemize}

\end{document}
